% Part: lambda-calculus
% Chapter: church-rosser
% Section: parallel-beta-eta-reduction

\documentclass[../../../include/open-logic-section]{subfiles}

\begin{document}

\olfileid{lam}{cr}{pbe}

\olsection{متوازی $\beta\eta$-تقلیل}

ایس حصے وچ اسیں متوازی $\beta\eta$-تقلیل لئی چرچ–روسر خاصیت
ثابت کردے آں۔ ایہ $\beta\eta$-تقلیل دے مطابق متوازی تقلیل اے۔

\begin{defn}[متوازی $\beta\eta$-تقلیل، $\beredpar$] \ollabel{defn:beredpar}
  ٹرماں اُتے \emph{متوازی $\beta\eta$-تقلیل} ($\beredpar$)
  استقرا نال ایویں متعین ہوندی اے۔
  \begin{enumerate}
    \item \ollabel{defn:beredpar1} $x \beredpar x$۔
    \item \ollabel{defn:beredpar2} جے $N \xrightarrow{\beta} N'$، تاں $\lambd[x][N] \beredpar
      \lambd[x][N']$۔
    \item \ollabel{defn:beredpar3} جے $P \beredpar P'$ تے $Q \beredpar Q'$، تاں $PQ \beredpar
      P'Q'$۔
    \item \ollabel{defn:beredpar4} جے $N \beredpar N'$ تے $Q \beredpar Q'$، تاں
      $(\lambd[x][N])Q \beredpar \Subst{N'}{Q'}{x}$۔
    \item \ollabel{defn:beredpar5} جے $N \beredpar N'$، تاں $\lambd[x][Nx]
      \beredpar N'$، بشرطیکہ $x \notin FV(N)$۔
  \end{enumerate}
\end{defn}

\begin{thm}\ollabel{thm:refl}
  $M \beredpar M$۔
\end{thm}

\begin{proof}
  مشق۔
\end{proof}

\begin{prob}
  \olref[lam][cr][pbe]{thm:refl} ثابت کرو۔
\end{prob}

\begin{defn}[$\beta\eta$-مکمل تقلیل]\ollabel{defn:becd}
  \emph{$\beta\eta$-مکمل تقلیل}~$\becd{M}$، جیہڑی~$M$ توں
  ملدی اے، ایویں متعین ہوندی اے۔
  \begin{align}
    \becd{x} &= x \ollabel{defn:becd1} \\
    \becd{(\lambd[x][N])} &= \lambd[x][\becd{N}] \ollabel{defn:becd2}\\
    \becd{(PQ)} &= \becd{P}\becd{Q} && \text{جے $P$، $\lambd$ تجرید نہیں}
    \ollabel{defn:becd3} \\
    \becd{((\lambd[x][N])Q)} &= \Subst{\becd{N}}{\becd{Q}}{x}
                             \ollabel{defn:becd4} \\
    \becd{(\lambd[x][Nx])} &= \becd{N} \ollabel{defn:becd5} & \text{جے $x
                                                      \notin FV(N)$}
  \end{align}
\end{defn}

\begin{lem}\ollabel{lem:comp}
  جے $M \beredpar M'$ تے $R \beredpar R'$، تاں $\Subst{M}{R}{y}
  \beredpar \Subst{M'}{R'}{y}$۔
\end{lem}

\begin{proof}
  $M \beredpar M'$ دے !!{derivation} اُتے استقرا نال۔

  پہلِیاں چار صورتاں عین \olref[pb]{lem:comp} والیاں نیں۔
  جے آخری قاعدہ \olref{defn:beredpar5} اے، تاں $M$ دی شکل
  $\lambd[x][Nx]$ تے $M'$ دی شکل $N'$ اے، کجھ $x$ تے~$N'$ لئی
  جتھے $x \notin
  FV(N)$ تے $N \beredpar N'$۔ اسیں دکھاؤنا اے کہ
  $\Subst{(\lambd[x][Nx])}{R}{y} \beredpar \Subst{N'}{R'}{y}$، یعنی
  $\lambd[x][\Subst{N}{R}{y} x] \beredpar \Subst{N'}{R'}{y}$۔ ایہ
  \olref{defn:beredpar}\olref{defn:beredpar5} تے استقرائی مفروضے توں ملدا اے۔
\end{proof}

\begin{lem}\ollabel{lem:cont}
  جے $M \beredpar M'$، تاں $M' \beredpar \becd{M}$۔
\end{lem}

\begin{proof}
  $M \beredpar M'$ دے !!{derivation} اُتے استقرا نال۔

  پہلِیاں چار صورتاں \olref[pb]{lem:cont} والیاں وانگ نیں۔
  جے آخری قاعدہ \olref{defn:beredpar5} اے، تاں $M$ دی شکل
  $\lambd[x][Nx]$ تے $M'$ دی شکل $N'$ اے، کجھ $x$، $N$، $N'$ لئی
  جتھے $x \notin FV(N)$ تے $N
  \beredpar N'$۔ اسیں دکھاؤنا اے کہ $N' \beredpar
  \becd{(\lambd[x][Nx])}$، یعنی $N' \beredpar \becd{N}$۔
  ایہ استقرائی مفروضے توں سیدھا ملدا اے۔
\end{proof}

\begin{thm}\ollabel{thm:cr}
  $\beredpar$ وچ چرچ–روسر خاصیت اے۔
\end{thm}

\begin{proof}
  \olref{lem:cont} توں سیدھا ملدا اے۔
\end{proof}

\end{document}
